% !PRINTED 593
\weberpage{593}
\begin{center}
\S.\,162.\\[4pt]
\gesperrt{Das Partial-Grundideal.}
\end{center}

Es möge jetzt
\weq{1}{\tau = t_1\omega_1 + t_2\omega_2 + \cdots + t_{mn}\omega_{mn}}
eine Basisform von $\mathfrak{o}$ sein. Dann wird es gewisse Substitutionen $\chi$ in $\Phi$ geben, und darunter immer die identische, die der Congruenz
\weq{2}{\tau \,|\, \chi \equiv \tau \pmod{\mathfrak{p}}}
genügen. Alle diese Substitutionen bilden eine in $\Phi$ enthaltene Gruppe $X$, die auch ein Theiler der Gruppe $\Psi$ ist, zu der $\mathfrak{p}$ gehört. Denn aus $\tau$ erhält man (nach \S.\,146) eine Basisform von $\mathfrak{p}$, wenn man für die Variablen $t$ gewisse lineare Functionen neuer Variablen mit ganzen rationalen Coëfficienten setzt; wenn also $\pi$ eine solche Basisform von $\mathfrak{p}$ ist, so folgt aus (2), dass $\pi \,|\, \chi$ durch $\pi$ theilbar und daher auch $\mathfrak{p} \,|\, \chi = \mathfrak{p}$ ist.

Nun genügt $\tau$ einer Gleichung $n^{\text{ten}}$ Grades $f(t) = 0$, deren Coëfficienten ganze Functionen in $R$ mit den Variablen $t_1, t_2 \ldots$ sind, und es ist, wenn $t$ eine neue Variable bedeutet, und $\tau_1, \tau_2, \ldots, \tau_n$ die in Bezug auf $R$ zu $\tau$ conjugirten Formen sind,
\weq{3}{f(t, t_1, t_2, \ldots) = f(t) = (t-\tau_1)(t-\tau_2) \ldots (t-\tau_n).}
Hieraus ergiebt sich nun, wenn wir wiederholt in die Potenz $p$ erheben, und auf die Zahlencoëfficienten von $f$ die Formel \S.\,161, (1) anwenden:
\weq{4}{f(t^P, t_1^P, t_2^P, \ldots) \equiv [f(t)]^P \equiv (t^P-\tau_1')(t^P-\tau_2') \ldots (t^P-\tau_n') \pmod{\mathfrak{p}},}
wenn $\tau_1', \tau_2', \ldots, \tau_n'$ dadurch aus $\tau_1, \tau_2, \ldots, \tau_n$ abgeleitet sind, dass die Variablen $t_1, t_2, \ldots$ durch $t_1^P, t_2^P, \ldots$ ersetzt sind.

Setzen wir nun $t = \tau$ in (4), so verschwindet $f(t)$ und es folgt, dass einer der Factoren des letzten Productes durch $\mathfrak{p}$ theilbar sein muss. Dies aber kann so ausgedrückt werden, dass es in $\Phi$ eine Substitution $\psi_0$ giebt, die der Bedingung
\weq{5}{\tau^P \equiv \tau' \,|\, \psi_0 \pmod{\mathfrak{p}}}
genügt.

Aus $\tau$ entstehen alle ganzen Zahlen in $\Omega$, wenn man für die Variablen ganze rationale Zahlen setzt. Wendet man dann
%CONT>>
